\section{Related Work}
\label{sec:related}

\paragraph{Vision-and-language navigation.}
Room-to-Room established instruction-following navigation with natural-language routes and path-based metrics~\cite{anderson2018vision}.
Most subsequent VLN settings assume egocentric observations and indoor or street-level motion.
AerialVLN extends instruction following to simulated UAV trajectories~\cite{liu2023aerialvln}, while CityNav provides real-city aerial trajectories and a geographic semantic map~\cite{lee2025citynav}.
Our setting instead uses north-aligned, georeferenced overhead imagery, metric actions, disaster-specific target semantics, and paired temporal observations.
These differences remove some 3D localization ambiguity but intensify small-object recognition and damage-assessment uncertainty.
We use hierarchical instruction decomposition and textual map summaries as supporting interfaces; the present experiments do not establish that these interfaces improve strict success rate.

\paragraph{Overhead disaster assessment.}
xBD, released with the xView2 challenge, provides large-scale pre/post-disaster satellite imagery and building-level damage labels~\cite{gupta2019xbd}.
It enables a navigation environment grounded in real disaster geometry rather than a purely synthetic simulator.
Winning xView2 entries commonly separate building localization from damage classification, often with U-Net-style encoders trained on pre-event imagery before a Siamese damage head.
We adopt only that localization idea as replaceable scaffolding: an event-disjoint ResNet34 U-Net proposes buildings, while damage probabilities still come from our dual-temporal classifier.
We do not load public competition weights, which would reintroduce event leakage through the proposer.
Object detectors such as YOLO~\cite{redmon2016you} and semantic segmentation architectures such as SegFormer~\cite{xie2021segformer} remain useful primitives, while modern vision-language models such as Qwen2-VL~\cite{wang2024qwen2vl} offer open-vocabulary interpretation.
However, detector confidence is not automatically a reliable probability of correctness, and overhead-to-aerial or cross-disaster transfer remains a major concern.

\paragraph{Calibration and uncertainty.}
Neural networks are often miscalibrated; temperature scaling is a simple validation-set procedure that can improve calibration while preserving the predicted class~\cite{guo2017calibration}.
We use expected calibration error (ECE), negative log-likelihood, and Brier score to audit probabilities.
Calibration alone does not guarantee good classification, distributional robustness, or useful actions.
Accordingly, we evaluate perception calibration separately from end-to-end reinspection effectiveness.

\paragraph{Active perception.}
Active perception treats camera motion as part of inference rather than as a passive data-collection step~\cite{bajcsy1988active}.
Receding-horizon next-best-view planning selects aerial viewpoints by expected exploration gain~\cite{bircher2016nbv}.
\ProjectName instantiates this idea with actions available to an overhead UAV: hold, descend, and recenter.
Our policy uses a bounded information-gain proxy rather than claiming an exact POMDP solution.
Its value must be measured against no-reinspection, random, and fixed-action controls at matched resource budgets.

\paragraph{Agentic earth observation and UAV systems.}
RemoteAgent studies intent understanding and capability-aware tool routing for
earth observation~\cite{yao2026remoteagent}; AerialClaw provides a modular
LLM--skill--runtime framework for autonomous aerial agents~\cite{li2026aerialclaw}.
\ProjectName builds on the broader closed-loop-agent view but asks a different,
narrow question: whether calibrated disaster-assessment uncertainty should
cause a physical reacquisition action.
It neither trains a general tool router nor treats framework breadth as the
scientific contribution.
